\section{Introduction}
\label{sec:intro}

An image of a tennis racket can make a ball plausible without making it
visible. A grounded vision--language model (VLM) should distinguish these
cases: it should recognize a clearly visible ball even without familiar
surrounding objects, and reject its presence when only those contextual cues
remain. This distinction matters for descriptions and visual question
answering, where a plausible but unsupported object can change the meaning
of a response. Object hallucination has been documented in image captioning
and modern VLMs~\cite{chair,pope}. Counterfactual removal and contextual
interventions further expose failures that aggregate accuracy can
conceal~\cite{removedobjects,causalhalbench}.

We focus on a narrow requirement: \emph{object-presence judgments should
respond to changes in the target object while remaining correct under
changes to surrounding context that preserve its presence}. These are
complementary requirements. A model that answers \emph{no} more often may
reduce false positives while missing more real objects. A model that repeats
the same answer across edited images may appear context-invariant while
being consistently wrong. Evaluation must therefore measure both sides of
the decision, and correctness across related images, rather than rewarding
invariance or lower hallucination in isolation.

Recent steering approaches offer alternatives to updating all model
parameters. VTI intervenes in visual and textual representations, GrAInS uses
gradient-based attribution, REVIS aims to restore suppressed visual
information, and RUDDER reinforces evidence extracted from residual
updates~\cite{vti,grains,revis,rudder}. These approaches motivate studying
where an intervention should act. However, their existence also raises the
standard for a new method: stronger visual features or sparse intervention
alone is not a sufficient methodological contribution.
HONES already connects causal head localization to sparse neuron
scaling~\cite{hones}, and recent work explicitly investigates visual-context
sensitivity and spurious associations~\cite{seeingknowing,purge}.

Our hypothesis concerns the \emph{selection criterion} and the
\emph{preservation-constrained intervention}. We ask whether heads selected
for a target-specific information-transfer effect provide more useful
intervention sites than equally many randomly chosen or high-magnitude
heads. Evidence-Preserving Head Steering (\method) first substitutes clean
head activations into target-occluded inputs and into matched background
controls. It ranks heads by the former effect after penalizing the latter,
then learns small head-wise gains on separate, balanced examples. A
preservation term discourages changes to the base model's next-token
distribution on examples it originally answers correctly.

Crucially, these steps have different meanings. Restoring a clean activation
to an object-occluded image is a diagnostic of information transfer; it can
restore an answer that is incorrect for the edited image. It does not show
that amplifying the same head improves deployment. Likewise, a matched
background occlusion controls some edit-related effects but does not isolate
semantic context. We therefore evaluate the learned gains using single
images, without donors, and reserve semantic-context interventions as a
separate test of the motivating hypothesis.

This work is organized around three proposed contributions: a
control-adjusted head-localization criterion; bounded gain fitting with
balanced supervision and preservation of correct base-model distributions;
and a comparison protocol that tests whether any improvement exceeds answer
bias calibration, arbitrary head interventions, and existing steering
methods. The first two have a CPU-tested implementation. Their empirical
benefit, the semantic-context protocol, and the adequacy of the comparison
to close prior work remain to be established. We do not claim a new
hallucination phenomenon or a new benchmark.

\begin{table}[t]
\centering
\small
\begin{tabular}{lll}
\toprule
Target object & Associated context & Correct answer \\
\midrule
Present & Present & Yes \\
Absent & Present & No \\
Present & Absent & Yes \\
Absent & Absent & No \\
\bottomrule
\end{tabular}
\caption{The intended diagnostic distinguishes target evidence from
associated context. Each question asks about the same target category.
These are experimental conditions, not observed model results. Correctness
must be assessed in every condition; constant answers are insufficient.}
\label{tab:conditions}
\end{table}
